\section{Introduction}
\label{sec:intro}

Multi-task learning (MTL) is a central pillar of modern artificial intelligence, enabling a single model to generalize across a diverse spectrum of domains and applications. While traditional multi-task training paradigms require simultaneous joint optimization over large-scale aggregated datasets, they suffer from extreme computational costs, data privacy concerns, and severe negative transfer (or task interference). To circumvent these bottlenecks, weight-space post-hoc model merging has emerged as a groundbreaking, training-free alternative. This paradigm allows separate expert models, independently fine-tuned from a shared pretrained initialization, to be directly composed into a single multi-task network without accessing the original training datasets.

Early model merging frameworks, such as \textbf{Task Arithmetic} \cite{task_vectors} and \textbf{TIES-Merging} \cite{ties_merging}, employ uniform, hand-tuned scaling coefficients to blend model weights. While simple and computationally lightweight, these static schemes ignore the complex, layer-specific correlations and conflicts between tasks, often resulting in sub-optimal joint performance. To resolve this, \textbf{Adaptive Model Merging (AdaMerging)} \cite{adamerging} introduced test-time adaptation (TTA), which dynamically learns layer-wise or task-wise merging coefficients by minimizing unsupervised entropy over a small calibration stream of unlabeled test queries.

\begin{figure}[t]
    \centering
    \includegraphics[width=0.9\linewidth]{comparison_plot.png}
    \caption{Model merging performance across post-training quantization (PTQ) schemas. Unregularized test-time adaptation (AdaMerging) collapses under low-precision formats. Our proposed \textbf{CR-PolySACM} framework combines global structural polynomial subspace constraints with local Clipping-Regularized Sharpness-Aware Minimization (CR-SACM), resolving the task-vector norm scale pathology and delivering robust performance across all target precisions (including setting a new state-of-the-art of \textbf{19.07\%} in INT4).}
    \label{fig:teaser}
\end{figure}

Despite the empirical success of TTA in continuous floating-point (FP32) spaces, we identify a critical, hitherto overlooked vulnerability in these adaptive paradigms: \textbf{Quantization-Operator Overfitting}. In real-world edge deployment, merged models are routinely quantized to low-precision formats (e.g., INT8 or INT4 symmetric or asymmetric schemas) to meet stringent hardware memory and latency constraints. We find that unregularized TTA optimizes the merging coefficients into extremely sharp local minima. While these sharp minima yield high performance in FP32, they are highly sensitive to even minor weight perturbations. When subjected to post-training quantization (PTQ), the resulting rounding noise triggers a catastrophic collapse in multi-task accuracy. This represents a fundamental hurdle for deploying test-time merged models on edge devices.

To investigate and resolve this vulnerability, we present a mathematically rigorous and empirical investigation into the limits of test-time flatness optimization and the power of structural subspace constraints:
\begin{enumerate}
    \item **The Task-Vector Norm Scale Pathology and CR-SACM:** We analyze unconstrained local flatness regularization, specifically \textbf{HessMerge} and its first-order minimax approximation \textbf{SACM}. Through a second-order Taylor expansion, we relate the quantization-induced loss gap directly to curvature. We identify a fundamental \emph{task-vector norm scale pathology} where a 50-fold discrepancy in layer-wise task-vector norms renders the optimizer blind to extremely sensitive, low-norm layers such as the final layer normalization layer (Layer 13). We propose \textbf{Clipping-Regularized SACM (CR-SACM)}, which clips task-vector norms to a robust minimum floor, successfully restoring optimizer sensitivity without triggering singular gradient explosion.
    \item **The Power of Subspace Constraints (PolyMerge):** We analyze subspace-constrained adaptive merging, where blending coefficients are restricted to a low-degree polynomial of network depth. By reducing the number of optimization parameters from 56 to 12, PolyMerge serves as a powerful global structural regularizer. This constraint prevents overfitting to the calibration stream, preserving the global, high-quality weight representations and naturally shielding the model against out-of-subspace noise.
    \item **The Integrated Subspace-Flatness Paradigm (CR-PolySACM):** We propose \textbf{CR-PolySACM} (referred to interchangeably as \textbf{PolySACM} throughout this work), which combines global structural depth-dependent polynomial subspace constraints with local CR-SACM flatness optimization. We demonstrate that CR-PolySACM successfully maps the optimization trajectory into flat, robust local minima within a highly regularized, well-conditioned search space, achieving exceptional performance gains.
\end{enumerate}

Our main contributions are summarized as follows:
\begin{itemize}
    \item We identify and analyze the phenomenon of \textbf{Quantization-Operator Overfitting} in test-time adaptive model merging, showing that unregularized coefficient optimization converges to sharp minima that collapse under post-training quantization.
    \item We define and analyze the \textbf{task-vector norm scale pathology} in weight-space sharpness optimization, showing that task-vector scale-blindness causes standard unnormalized flatness regularizers to collapse or fail.
    \item We introduce \textbf{Clipping-Regularized SACM (CR-SACM)}, a robust and mathematically sound perturbation mechanism that balances scale sensitivity across layer groups and resolves scale-blindness without triggering gradient explosion.
    \item We propose and evaluate \textbf{CR-PolySACM}, a unified model merging framework that combines structured depth-dependent polynomial subspace constraints with local CR-SACM flatness optimization to achieve stable, high-performance composition.
    \item We conduct a comprehensive multi-schema evaluation (including FP32, INT8 symmetric/asymmetric tensor-wise and channel-wise, and INT4 symmetric channel-wise quantization) across four diverse visual domain datasets using a Vision Transformer backbone, demonstrating that CR-PolySACM consistently outperforms standard PolyMerge and sets a new state of the art of \textbf{19.07\%} in INT4.
\end{itemize}
